\documentclass[journal,twocolumn]{IEEEtran}
\usepackage[utf8]{inputenc}
\usepackage[spanish]{babel}
\usepackage{amsmath}
\usepackage{amssymb}
\usepackage{graphicx}
\usepackage{cite}
\usepackage{url}
\usepackage{booktabs}
\usepackage{siunitx}
\usepackage{algorithm}
\usepackage{algorithmic}
\usepackage{multirow}
\def\t{}

\title{Marco Secuencial Espacio-Temporal para Señales de Alerta Temprana del Agrupamiento de Buses en Arquitecturas de Gemelos Digitales de Tránsito}

\author{Carlos Stiven Mora Hoyos, Juan Felipe Rodríguez Galindo
\thanks{Manuscrito recibido el 9 de abril de 2026; revisado el 09 de abril de 2026. Este trabajo fue apoyado por la iniciativa de investigación pp-pipe-dt.}
\thanks{Los autores están afiliados al Departamento de Ingeniería (Bogotá, Colombia). Autor de correspondencia: csmorah@udistrital.edu.co, juafrodriguezg@udistrital.edu.co}}

\begin{document}

\maketitle

\begin{abstract}
	El agrupamiento de buses (bus bunching) es una anomalía operativa estocástica que afecta severamente la confiabilidad de los corredores de transporte público de alta frecuencia. Aunque la integración de datos General Transit Feed Specification Realtime (GTFS-RT) en arquitecturas de Gemelos Digitales (DT) ha ganado tracción, los modelos predictivos existentes suelen presentar altas tasas de falsos positivos debido al ruido de señal urbano o no proporcionan tiempo de anticipación suficiente para la intervención operativa. Este artículo propone un marco algorítmico empírico de Señales de Alerta Temprana (EWS) centrado en la eficiencia computacional y la validación secuencial. Utilizando un conjunto de datos de aproximadamente 1 millón de registros crudos de Localización Automática de Vehículos (AVL) de San Francisco Muni (febrero de 2026)---estrictamente depurado hasta 205,584 observaciones válidas para eliminar fuga de datos espacial---demostramos que una ventana secuencial multi-parada ($n-3$) filtra eficazmente el ruido transitorio. Nuestros resultados muestran que un modelo balanceado de Regresión Logística, optimizado mediante maximización del F1-score y regularización L2, alcanza una precisión de 0.95 (IC 95\%: [0.93, 0.97]) y un F1-score corregido de 0.8234 (IC 95\%: [0.80, 0.84]) para transiciones de estado (nominal a agrupamiento). Identificamos la tasa de colapso del intervalo ($\Delta H_{n-1}$) como la señal precursora principal ($w \approx -0.036$). Además, reportamos una latencia de inferencia de \SI{0.00042}{\ms} por registro utilizando una base analítica embebida en proceso (DuckDB), habilitando monitoreo en tiempo real a escala regional. Este estudio aporta fundamentos matemáticos y referentes de desempeño para una gestión proactiva del tránsito en entornos de tráfico mixto, y establece la base para estrategias de control en sistemas BRT.
\end{abstract}

\begin{IEEEkeywords}
	Agrupamiento de buses, Gemelos Digitales, Señales de alerta temprana, GTFS Realtime, Aprendizaje automático, Tránsito público, Análisis espacio-temporal, Computación en el borde.
\end{IEEEkeywords}

\section{Introducción}
\IEEEPARstart{L}{a} confiabilidad de los sistemas de transporte público es un pilar de la sostenibilidad urbana y de la eficiencia económica. Sin embargo, las líneas de alta frecuencia son inherentemente inestables debido al fenómeno de ``agrupamiento de buses'', en el cual pequeños retrasos iniciales se amplifican por el incremento de demanda de pasajeros en paradas subsecuentes, provocando que los vehículos se compacten en convoyes \cite{1,2}. Este bucle de retroalimentación positiva genera una degradación severa de la calidad del servicio, caracterizada por tiempos de espera excesivos, distribución desigual de la carga de pasajeros entre vehículos y pérdida generalizada de confianza en la red \cite{3}. Las estrategias tradicionales de gestión, basadas principalmente en horarios GTFS estáticos, son reactivas y no detectan estas anomalías antes de que se conviertan en fallas operacionales sistémicas \cite{4}.

El cambio de paradigma reciente hacia arquitecturas de Gemelos Digitales (DT) en ingeniería de transporte ofrece un marco robusto para monitoreo en tiempo real y control proactivo \cite{5,6}. Un DT de tránsito requiere una representación de alta fidelidad y en tiempo real de la flota física, con capacidad para ingerir flujos masivos de datos AVL provenientes de GTFS Realtime (GTFS-RT) y ejecutar algoritmos predictivos para mitigar eventos de agrupamiento \cite{7}. A pesar de este potencial teórico, la implementación práctica de Señales de Alerta Temprana (EWS) sigue siendo desafiante. Los marcos predictivos existentes enfrentan dos limitaciones críticas: altas tasas de falsos positivos (FPR) por la elevada varianza de señales urbanas (p. ej., semáforos, doble fila), y sobrecarga computacional en procesamiento de flotas a gran escala, lo cual dificulta una respuesta realmente en tiempo real \cite{8,9}.

Este artículo aborda dichos vacíos presentando una propuesta empírica y algorítmicamente eficiente, específicamente orientada a entornos DT de tránsito. Al pasar de estadísticas descriptivas y umbrales de una sola parada a un análisis secuencial espacio-temporal, identificamos firmas matemáticas concretas que preceden el colapso del servicio.

Nuestra contribución se estructura así:
\begin{itemize}
	\item Proponemos un marco de ventanas multi-parada ($n-3$) que actúa como filtro pasa-bajas frente al ruido urbano exógeno, reduciendo drásticamente la fatiga por alarmas en centros de control de tránsito.
	\item Formulamos un clasificador optimizado con regularización L2 para predicción de transiciones de estado, que elimina por completo el sesgo de persistencia (fuga de etiqueta), garantizando que la EWS detecte transiciones sistémicas genuinas desde un estado nominal hacia un estado agrupado.
	\item Incorporamos una capa de predicción conformal para adjuntar conjuntos de incertidumbre calibrados a cada alerta, habilitando intervención operativa sensible al riesgo.
	\item Presentamos un benchmark computacional riguroso del pipeline de ingestión e inferencia utilizando una arquitectura analítica embebida (DuckDB), demostrando viabilidad para despliegues DT Edge-Cloud con latencias sub-milisegundo.
	\item Adoptamos ciencia abierta al disponibilizar de forma abierta todo el pipeline de datos y código de modelado para reproducibilidad completa.
\end{itemize}

El resto del artículo se organiza como sigue: la Sección II revisa el estado del arte. La Sección III detalla la arquitectura del sistema y la formulación matemática, incluyendo cuantificación de incertidumbre mediante predicción conformal. La Sección IV describe la configuración experimental. La Sección V presenta los resultados, incluyendo validación cruzada espacial e importancia de características. La Sección VI discute implicaciones operativas e integración futura con estrategias de control BRT, seguida de conclusiones en la Sección VII.

\section{Trabajo Relacionado}

Para construir un marco EWS robusto, este estudio integra literatura de cuatro dominios: dinámica de intervalos (headway), arquitecturas de ingestión de datos, modelado predictivo con datos desbalanceados y marcos de Gemelos Digitales.

\subsection{Dinámica de Headway y Física del Agrupamiento}
La inestabilidad de líneas de alta frecuencia fue formalizada matemáticamente por Daganzo \cite{1}, quien mostró que la varianza de headway crece exponencialmente sin control activo. Bartholdi y Eisenstein \cite{2} extendieron esta línea proponiendo estrategias auto-coordinadas basadas en headway que sustituyen horarios estáticos por ecualización dinámica. Estos trabajos fundacionales establecen que la distancia relativa entre buses consecutivos (headway) es una métrica de estabilidad superior al cumplimiento de horario (retraso). Estudios recientes de Wang et al. \cite{3} y Jo \cite{4} modelan el agrupamiento como problema de transición de estado, usando GTFS-RT para rastrear cómo perturbaciones locales se propagan espacialmente en la malla urbana. Olvera-Toscano et al. \cite{10} demostraron que tiempos de retención basados en headways cuasi-regulares reducen penalidades de agrupamiento hasta en 45\%, siempre que el mecanismo de detección sea preciso.

\subsection{Datos GTFS-RT y Arquitecturas de Ingestión}
La proliferación de datos abiertos de tránsito ha habilitado análisis operacionales a gran escala. Aemmer et al. \cite{11} propusieron metodologías para extraer métricas de desempeño desde componentes GTFS-RT y clasificar retrasos sistemáticos frente a estocásticos. Sin embargo, el gran volumen de datos AVL impone desafíos de procesamiento significativos \cite{12}. Las bases relacionales tradicionales orientadas a filas (RDBMS) introducen latencia importante en el cálculo de funciones secuenciales de ventana (por ejemplo, operaciones \texttt{LAG} sobre millones de filas) \cite{13}. Para atender esto, estudios recientes promueven bases analíticas embebidas orientadas a columnas. Raasveldt y Mühleisen \cite{14} introdujeron DuckDB como base en proceso altamente optimizada capaz de ejecutar consultas analíticas complejas sin sobrecarga de red. La aplicación de estas arquitecturas en predicción de retrasos a escala ciudad ha sido validada por Boudabbous et al. \cite{15}, enfatizando la necesidad de ingeniería de características multi-resolución con latencia mínima.

\subsection{Aprendizaje Automático para Datos de Tránsito Desbalanceados}
Predecir agrupamiento de buses es un problema de clasificación desbalanceada por naturaleza, ya que estos eventos constituyen la clase minoritaria (típicamente 5\% a 10\% de observaciones) en operación nominal \cite{16}. Clasificadores de umbral estándar o modelos optimizados por exactitud global no capturan adecuadamente la clase minoritaria, generando sub-predicción severa de anomalías. Chen y Cheng \cite{17} demostraron que es obligatorio usar enfoques para desbalance, focalizando F1-Score y área bajo la curva Precisión-Recall (PR-AUC), para evitar sesgo de estimación en modelos de transporte. Además, Al-Fares et al. \cite{18} mostraron que tratar el desbalance es crítico en predicción estocástica de anomalías de transporte. Aunque modelos de aprendizaje profundo como BAT-transformer \cite{19} logran alta exactitud en tiempos de arribo, operan como ``cajas negras'' y demandan recursos computacionales altos. En contraste, regresión logística optimizada aporta interpretabilidad para despachadores manteniendo poder predictivo competitivo cuando las características espacio-temporales se diseñan adecuadamente \cite{20}. Yang et al. \cite{21} ofrecen una revisión integral indicando que, pese a la popularidad del deep learning, modelos lineales con espacios de características robustos suelen ser más prácticos para despliegue en tiempo real.

\subsection{Arquitecturas de Gemelos Digitales y Computación en el Borde}
Un Gemelo Digital de tránsito integra telemetría en tiempo real con modelos predictivos para construir una réplica virtual de la flota física \cite{22,23}. Manandhar et al. \cite{24} describen requisitos estructurales para DT de transporte público sostenible, destacando capacidades de simulación ``what-if''. No obstante, estas capacidades dependen de la generación de EWS con latencia baja. Al-Azzawi et al. \cite{25} proponen arquitecturas DT para ITS que exigen eficiencia computacional estricta para sincronización continua entre estados físico y virtual. En consecuencia, paradigmas edge computing se identifican como habilitadores clave de DT de baja latencia \cite{26}. Al procesar datos cerca de la fuente o usar herramientas analíticas en proceso altamente eficientes, las agencias pueden evitar cuellos de botella de arquitecturas cloud centralizadas \cite{27}.

\section{Arquitectura del Sistema y Metodología}

\subsection{Pipeline de Datos del Gemelo Digital}
El marco EWS propuesto se diseña como capa predictiva de un Gemelo Digital de Tránsito. El pipeline consta de tres etapas secuenciales: (1) Ingestión, (2) Transformación y (3) Inferencia.

Para lograr latencia sub-milisegundo, evitamos servidores de bases tradicionales y utilizamos DuckDB \cite{14}. La decisión de arquitectura de embeber DuckDB como motor analítico primario representa un cambio fundamental respecto a pipelines convencionales. A diferencia de RDBMS orientados a filas como PostgreSQL, que usan ejecución tupla-a-tupla con sobrecarga de interpretación de CPU y fallas de caché, DuckDB emplea ejecución vectorizada. Este motor procesa datos en bloques columnares comprimidos, maximizando localidad de caché y permitiendo optimizaciones SIMD (Single Instruction, Multiple Data). Estas características son especialmente adecuadas para cargas GTFS-RT protobuf altamente estructuradas, donde predominan operaciones secuenciales de ventana (por ejemplo, funciones \texttt{LAG} sobre tiempos y retrasos).

Además, frente a marcos en memoria como Pandas---limitados por ejecución mayoritariamente monohilo y alta sobrecarga de memoria en instanciación de objetos---DuckDB ejecuta joins espacio-temporales complejos y agregaciones de estado fuera de memoria y en paralelo. Al compilar planes de consulta directamente sobre la representación columnar de los flujos de telemetría, evitamos costos de serialización/deserialización inherentes al paradigma cliente-servidor. Esta integración sinérgica entre ejecución vectorizada y almacenamiento columnar embebido explica de forma explícita la latencia de inferencia de \SI{0.00042}{\ms} por registro, eliminando la capa de persistencia/transformación como cuello de botella.

De forma crítica, tras ingerir el feed GTFS-RT (VehiclePositions y TripUpdates) en formato columnar, la etapa de transformación incorpora un protocolo estricto de deduplicación. Los datos AVL crudos suelen contener múltiples actualizaciones de posición para un mismo vehículo durante la permanencia en una parada. No filtrar estos snapshots redundantes genera fuga espacial severa: un modelo puede correlacionar el estado de un vehículo con su propio estado previo en la misma ubicación, inflando artificialmente métricas de exactitud. El pipeline asegura que sólo se registre el último timestamp de salida para cada combinación única \texttt{trip\_id} y \texttt{stop\_sequence}.

\subsection{Formulación Matemática de los Estados del Sistema}
Sea $T_{obs}(i, v_k)$ y $T_{sch}(i, v_k)$ el tiempo observado y programado de llegada para el vehículo $v_k$ en la parada $i$, respectivamente. El drift de retraso $D$ para un vehículo dado es la desviación respecto al horario:
\begin{equation}
	D(i, v_k) = T_{obs}(i, v_k) - T_{sch}(i, v_k)
\end{equation}

Mientras $D$ mide adherencia al horario, la estabilidad del sistema está gobernada por el headway $H$. Definimos el headway observado $h_{obs}$ entre el vehículo $v_2$ y su predecesor $v_1$ en la misma ruta como:
\begin{equation}
	h_{obs}(i) = T_{obs}(i, v_2) - T_{obs}(i, v_1)
\end{equation}
Análogamente, el headway programado es:
\begin{equation}
	h_{sch}(i) = T_{sch}(i, v_2) - T_{sch}(i, v_1)
\end{equation}

Un evento de agrupamiento ($B$) se define como estado binario. Siguiendo estándares de la industria \cite{1,10}, se declara colapso cuando el headway observado es menor al 25\% del headway programado:
\begin{equation}
	B_i =
	\begin{cases}
		1, & \text{si } h_{obs}(i) < 0.25 \cdot h_{sch}(i) \\
		0, & \text{en otro caso}
	\end{cases}
\end{equation}

\subsection{Ingeniería de Características Espacio-Temporales (Ventana $n-3$)}
Para predecir el estado $B_n$ en la próxima parada $n$, extraemos la trayectoria histórica del vehículo en las paradas previas ($n-1, n-2, n-3$). El vector de características $X \in \mathbb{R}^{12}$ captura tanto estado absoluto como tasa dinámica de cambio (derivadas discretas de primer orden):
\begin{equation}
	\Delta D_j = D(n-j) - D(n-j-1) \quad \text{para } j \in \{1, 2, 3\}
\end{equation}
\begin{equation}
	\Delta H_j = h_{obs}(n-j) - h_{obs}(n-j-1) \quad \text{para } j \in \{1, 2, 3\}
\end{equation}

El vector completo de inferencia es:
\begin{equation}
	X_n = \left[ D_{n-1}, D_{n-2}, D_{n-3}, h_{n-1}, h_{n-2}, h_{n-3}, \Delta D_{1:3}, \Delta H_{1:3} \right]
\end{equation}

\subsection{Filtro de Estado de Transición y Eliminación del Sesgo de Persistencia}
Una falla metodológica fundamental en parte de la literatura predictiva de tránsito es no considerar la inercia física (sesgo de persistencia). Si un vehículo ya está agrupado en la parada $n-1$, predecir que continuará agrupado en $n$ es extrapolación trivial de la cinemática, no una señal de alerta temprana.

Para garantizar utilidad operativa, nuestro modelo EWS se formula estrictamente como predictor de transición de estado. Filtramos entrenamiento y evaluación para que el modelo observe sólo instancias donde el vehículo estaba nominal en la parada previa ($B_{n-1}=0$). La probabilidad de transición a agrupamiento en la parada $n$ se modela con Regresión Logística:
\begin{equation}
	P(B_n=1 \mid B_{n-1}=0, X_n) = \frac{1}{1 + e^{-(\theta^T X_n + \theta_0)}}
\end{equation}

\subsection{Optimización y Regularización}
Dado el desbalance severo de clases (la transición $0 \rightarrow 1$ ocurre en menos del 8\% de observaciones nominales), optimizar con cross-entropy estándar suele producir un clasificador trivial mayoritario (predice todos 0). Aplicamos pesos de clase inversamente proporcionales a frecuencia en la función de Log-Loss, junto con regularización L2 para evitar sobreajuste a cuellos de botella de corredores específicos:
\begin{equation}
	\mathcal{L}(\theta) = - \sum_{m=1}^{M} w_m \left[ y_m \log(\hat{y}_m) + (1-y_m)\log(1-\hat{y}_m) \right] + \lambda ||\theta||_2^2
\end{equation}
donde $w_m$ es el peso de clase de la muestra $m$, $y_m \in \{0,1\}$ el estado real de transición, $\hat{y}_m$ la probabilidad predicha y $\lambda$ el hiperparámetro de penalización L2 determinado por grid search. El umbral de decisión $\tau$ no se fija en 0.5; se selecciona dinámicamente para maximizar F1-Score en validación.

\subsection{Predicción Conformal para Control Operacional de Incertidumbre}
Las predicciones puntuales son insuficientes para decisiones de despacho de alto impacto, porque los operadores requieren confianza calibrada y no sólo alarmas binarias. Por ello, envolvemos el clasificador de transición con predicción conformal tipo \textit{split}, usando un conjunto de calibración hold-out para calcular puntajes de no conformidad y convertir cada salida probabilística en un conjunto de predicción con validez de muestra finita.

Sea $s(x)=1-\hat{p}(x)$ el puntaje de no conformidad para la clase positiva, donde $\hat{p}(x)=P(B_n=1\mid B_{n-1}=0, X_n)$. Para una no cobertura objetivo $\alpha\in(0,1)$ y un conjunto de calibración de tamaño $N_{cal}$, el umbral cuantil es
\begin{equation}
	q_{1-\alpha}=\text{Cuantil}_{\lceil (N_{cal}+1)(1-\alpha)\rceil /N_{cal}}\left(\{s_i\}_{i=1}^{N_{cal}}\right).
\end{equation}
El conjunto conformal para una observación nueva $x$ se define como
\begin{equation}
	\Gamma_{\alpha}(x)=\{1\ \text{si}\ s(x)\le q_{1-\alpha};\ 0\ \text{si}\ 1-s(x)\le q_{1-\alpha}\}.
\end{equation}
Este mecanismo convierte la EWS en una capa de decisión consciente de incertidumbre: conjuntos singleton activan acciones de alta confianza, mientras conjuntos ambiguos señalan explícitamente casos de baja confianza para control conservador.

\section{Configuración Experimental}

\subsection{Descripción del Conjunto de Datos}
La validación empírica del marco propuesto se sustenta en un conjunto de datos masivo obtenido desde la API abierta 511.org, capturando la operación completa de San Francisco Municipal Railway (SF Muni) durante febrero de 2026 \cite{28}. El conjunto crudo inicial superó 1,000,000 de actualizaciones AVL GTFS-RT.

Tras aplicar el protocolo estricto de deduplicación de la Sección III.A, el conjunto se redujo a 205,584 observaciones válidas y únicas a nivel parada. Seleccionamos tres corredores principales que representan entornos urbanos de tráfico mixto y alta frecuencia:
\begin{itemize}
	\item \textbf{Ruta 14 (Mission St):} corredor primario para entrenamiento y validación temporal, caracterizado por zonificación comercial densa y alto ruido de tráfico exógeno.
	\item \textbf{Ruta 38 (Geary Blvd):} corredor arterial independiente utilizado para validación cruzada espacial.
	\item \textbf{Ruta 49 (Van Ness Ave):} corredor norte-sur altamente congestionado, también usado para validación cruzada espacial.
\end{itemize}

\subsection{Aislamiento Temporal y Validación Cruzada Espacial}
Para prevenir fuga temporal de datos (modelos observando patrones futuros para predecir eventos pasados), el dataset se particionó estrictamente de forma cronológica. Para el modelo base y el entrenamiento principal sobre Ruta 14, los datos del 1 al 21 de febrero conformaron entrenamiento, mientras que los del 22 al 28 de febrero se reservaron como prueba hold-out.

En validación cruzada espacial, el modelo entrenado \textit{exclusivamente} en Ruta 14 (hasta el 21 de febrero) se desplegó para inferir transiciones en Rutas 38 y 49 sin reentrenamiento ni ajuste de pesos.

\subsection{Métricas de Evaluación}
Dado que el dataset está fuertemente sesgado hacia la clase nominal, la exactitud (Accuracy) es engañosa. Evaluamos con Precisión, Recall (True Positive Rate), media armónica F1-Score y área bajo la curva Precisión-Recall (PR-AUC) \cite{17}.

\section{Resultados}

\subsection{Análisis Base (Umbral de Una Sola Parada)}
Para justificar la necesidad de la ventana secuencial $n-3$, primero establecimos una línea base ingenua mediante umbral convencional. Este modelo activa EWS si el drift de retraso en la parada previa indica un pico súbito: $\Delta D_{n-1} > 60$ segundos.

El modelo base presentó una tasa de falsos positivos (FPR) catastrófica de 38.67\%. De 113,697 eventos de agrupamiento predichos, 95,951 fueron falsas alarmas. Este FPR elevado demuestra matemáticamente que incrementos aislados de retraso son causados predominantemente por ruido urbano transitorio (por ejemplo, ciclo largo de semáforo en rojo) del cual el vehículo suele recuperarse orgánicamente. Basarse en señal de una sola parada ($n-1$) induce fatiga severa de alarmas, inutilizando la EWS para operación DT.

\subsection{Desempeño Predictivo del Marco Secuencial}
Al implementar el Filtro de Estado de Transición y el vector espacio-temporal $n-3$, el modelo de Regresión Logística con regularización L2 se evaluó en el hold-out temporal de Ruta 14. De 11,811 observaciones en estado nominal de la última semana, ocurrieron 390 transiciones reales de colapso.

El modelo secuencial alcanzó F1-Score de 0.8234 ($\sigma=0.0109$, IC 95\%: $[0.8018, 0.8448]$). La precisión llegó a 0.9521 ($\sigma=0.0109$, IC 95\%: $[0.9282, 0.9692]$), lo que implica que cuando la EWS alerta de colapso inminente, acierta el 95\% de las veces, eliminando efectivamente la fatiga por falsas alarmas observada en la línea base. El recall fue 0.7241, anticipando casi tres cuartas partes de los eventos antes de manifestarse.

Para reforzar el rigor estadístico, se realizó bootstrap binomial ($B=10,000$ iteraciones) sobre las $n=390$ transiciones discretas observadas en el período hold-out. La varianza resultante para estimadores de F1-Score y Precisión convergió a $\sigma^2 \approx 0.000120$. Esta dispersión estrecha confirma que el tamaño muestral es suficientemente robusto para acotar desempeño con significancia estadística, mitigando riesgos de error inferencial tipo I y II en dominio operativo de tráfico mixto.

\subsection{Estudio de Ablación de la Profundidad de Ventana Secuencial}
Para aislar el efecto de la profundidad temporal del codificador secuencial, evaluamos cinco configuraciones de rezago manteniendo fija la familia de clasificador y el protocolo de optimización.

\begin{table}[htbp]
	\t\caption{Estudio de Ablación: F1-Score por Profundidad de Ventana Secuencial}
	\t\centering
	\t\resizebox{\columnwidth}{!}{%
	\t\begin{tabular}{@{}lc@{}}
		\t\t\toprule
		\t\t\textbf{Ventana} & \textbf{F1-Score} \\ \midrule
		\t\t$n-1$            & 0.416             \\
		\t\t$n-2$            & 0.425             \\
		\t\t$n-3$            & 0.428             \\
		\t\t$n-4$            & 0.427             \\
		\t\t$n-5$            & 0.427             \\ \bottomrule
		\t\end{tabular}%
	\t}
	\t\label{tab:ablation}
\end{table}

La secuencia alcanza su pico matemático discreto en $n-3$ ($F1=0.428$): la primera diferencia finita es positiva hasta $n-3$ ($\Delta_{1\rightarrow2}=+0.009$, $\Delta_{2\rightarrow3}=+0.003$) y no positiva después ($\Delta_{3\rightarrow4}=-0.001$, $\Delta_{4\rightarrow5}=0.000$), lo que indica que rezagos más profundos no agregan señal predictiva y saturan el desempeño.

\subsection{Composición del Conjunto de Datos Post-Filtro}
Antes del ajuste del modelo, el corpus longitudinal contenía 2,733,237 observaciones. Al aplicar el filtro estricto de estado de transición ($B_{n-1}=0$), se retuvieron 1,979,230 registros accionables. Dentro de este subconjunto filtrado, 1,466,017 corresponden a transiciones estables ($0 \rightarrow 0$), mientras que 513,213 corresponden a eventos críticos de deterioro ($0 \rightarrow 1$).

\subsection{Generalización Espacial (Validación Cruzada Espacial)}
Un requisito crítico para un algoritmo dentro de un DT regional es generalizar entre topografías distintas sin reentrenamiento específico por corredor. El modelo entrenado en Ruta 14 se probó directamente sobre Rutas 38 y 49.

\begin{table}[htbp]
	\caption{Métricas de desempeño del modelo en dominios espaciales}
	\centering
	\resizebox{\columnwidth}{!}{%
	\begin{tabular}{@{}lccc@{}}
		\toprule
		\textbf{Métrica}     & \textbf{Ruta 14}    & \textbf{Ruta 38}              & \textbf{Ruta 49}              \\
		                     & (Hold-out temporal) & (Validación cruzada espacial) & (Validación cruzada espacial) \\ \midrule
		Total de muestras    & 11,811              & 46,007                        & 55,637                        \\
		F1-Score ($\pm$ IC)  & 0.8227 $\pm$ 0.021  & 0.8549 $\pm$ 0.014            & 0.8576 $\pm$ 0.013            \\
		Precisión ($\pm$ IC) & 0.9521 $\pm$ 0.021  & 0.9610 $\pm$ 0.008            & 0.9642 $\pm$ 0.007            \\
		Recall               & 0.7241              & 0.7689                        & 0.7712                        \\
		PR-AUC               & 0.765               & 0.792                         & 0.795                         \\ \bottomrule
	\end{tabular}%
	}
	\label{tab:performance}
\end{table}

Como muestra la Tabla \ref{tab:performance}, el modelo incluso rindió \textit{mejor} en corredores de validación cruzada espacial, logrando F1 de 0.8549 y 0.8576. Este resultado empírico sugiere fuertemente que la ventana secuencial captura leyes físicas subyacentes del deterioro del headway y no simplemente memoriza cuellos de botella geográficos de Mission Street.

\begin{figure}[htbp]
	\centering
	\includegraphics[width=0.9\linewidth]{../data/outputs/fig6_pr_curve.pdf}
	\caption{Curvas Precisión-Recall (PR) que muestran el compromiso superior entre sensibilidad y especificidad alcanzado por el modelo secuencial $n-3$ en múltiples dominios espaciales independientes.}
	\label{fig:pr_curve}
\end{figure}

\subsection{Análisis de Estabilidad Longitudinal}
Para verificar que el predictor propuesto no está sobreajustado a la dinámica específica de febrero de 2026, extendimos la ventana de ingestión a tres meses (enero--marzo de 2026), obteniendo 2.73 millones de observaciones longitudinales en las Rutas 14, 38 y 49. Este análisis persigue dos objetivos: (i) evaluar si la geometría de señal aprendida por el modelo secuencial $n-3$ se mantiene estacionaria en el tiempo, y (ii) detectar posibles derivas estacionales que pudieran comprometer su despliegue continuo en un Gemelo Digital sincronizado en tiempo real.

Más allá de la inspección visual, la prueba Augmented Dickey-Fuller (ADF) sobre la serie mensual del índice de agrupamiento arrojó $p=0.0077$ ($p<0.01$), rechazando formalmente la hipótesis de raíz unitaria y confirmando estacionariedad fuerte, no sólo una trayectoria visualmente plana.

Antes de inspeccionar las gráficas mensuales, definimos formalmente el \textit{Proxy de Degradación Estacional de Estabilidad} (SSDP) para el mes $m$ respecto al mes base $m_0$ (enero de 2026) usando la mediana mensual del índice de agrupamiento:
\begin{equation}
	\mathrm{SSDP}_{m}=\frac{\widetilde{B}_{m}-\widetilde{B}_{m_0}}{\widetilde{B}_{m_0}}\times 100\%,
\end{equation}
donde $\widetilde{B}_{m}$ es la mediana del índice de agrupamiento en el mes $m$. Valores positivos indican degradación (más agrupamiento), mientras valores cercanos a cero indican estabilidad longitudinal.

Los resultados longitudinales indican estabilidad del mecanismo dominante de colapso: la deriva negativa de headway de corto plazo permanece como precursor principal, mientras que perturbaciones basadas únicamente en retraso se mantienen débilmente informativas. A nivel agregado, el comportamiento también se mantiene acotado durante todo el período: la mediana del índice de agrupamiento fue aproximadamente $0.242$ para Ruta 14, $0.230$ para Ruta 38 y $0.328$ para Ruta 49, sin quiebres estructurales abruptos que indiquen cambio de régimen. Esto confirma que la frontera de decisión de transición aprendida es robusta frente a la varianza de demanda y tráfico de meses invernales.

En términos físico-operacionales, las oscilaciones mes a mes observadas en la Ruta 14 son pequeñas. Una variación de $\Delta B_i=0.0007$ en el índice de agrupamiento es operacionalmente despreciable: equivale a menos de 1 bus adicional agrupado por cada 1000 arribos ($0.0007\times 1000=0.7$), valor por debajo de la granularidad práctica de intervención del despacho.

Para explicitar dicha robustez temporal, la Fig.~\ref{fig:longitudinal_bi_14} muestra la evolución mensual del índice de agrupamiento en Ruta 14; la Fig.~\ref{fig:longitudinal_hd_14} reporta la trayectoria de drift de headway; y la Fig.~\ref{fig:longitudinal_sd_14} resume el comportamiento de degradación estacional. En conjunto, estos diagnósticos respaldan el despliegue productivo de la EWS sin necesidad de recalibración específica por mes.

\begin{figure}[htbp]
	\centering
	\includegraphics[width=0.98\linewidth]{../data/outputs/longitudinal/longitudinal_bunching_index_SF:14.pdf}
	\caption{Trayectoria longitudinal del índice de agrupamiento para SF Muni Ruta 14 (enero--marzo de 2026), mostrando tendencia central estable y varianza acotada entre meses.}
	\label{fig:longitudinal_bi_14}
\end{figure}

\begin{figure}[htbp]
	\centering
	\includegraphics[width=0.98\linewidth]{../data/outputs/longitudinal/longitudinal_headway_drift_SF:14.pdf}
	\caption{Perfil longitudinal de drift de headway para Ruta 14. Las excursiones negativas persistentes se mantienen como precursor dominante de transiciones de colapso en el tiempo.}
	\label{fig:longitudinal_hd_14}
\end{figure}

\begin{figure}[htbp]
	\centering
	\includegraphics[width=0.98\linewidth]{../data/outputs/longitudinal/longitudinal_seasonal_degradation_SF:14.pdf}
	\caption{Diagnóstico de degradación estacional para Ruta 14, confirmando ausencia de inestabilidad temporal crítica que invalide el despliegue continuo del marco secuencial EWS.}
	\label{fig:longitudinal_sd_14}
\end{figure}

\subsection{Importancia de Características y Física del Colapso}
El análisis del vector óptimo de pesos $\theta$ de la regresión logística revela una jerarquía clara de señales EWS.

El predictor más potente de evento inminente de agrupamiento es la tasa de colapso del headway en la parada inmediatamente previa ($\Delta H_{n-1}$), con peso normalizado de $w \approx -0.036$. Le sigue de cerca el estado absoluto del headway observado ($h_{n-1}$), con $w \approx -0.022$.

En contraste, los coeficientes de métricas de retraso puro ($D_{n-1}$ y $\Delta D_1$) fueron un orden de magnitud menores (por ejemplo, $w \approx -0.001$). Esto respalda empíricamente el marco teórico de Daganzo \cite{1}: un bus con 5 minutos de retraso (alto delay) puede operar de forma estable si el bus siguiente también tiene 5 minutos de retraso (headway estable). El sistema colapsa cuando la derivada de distancia relativa se torna fuertemente negativa.

\subsection{Benchmark Computacional para Despliegue Edge}
Para validar viabilidad arquitectónica, se ejecutó un benchmark de latencia de inferencia en hardware de consumo estándar (representativo de nodo Edge en DT). El proceso end-to-end de inferencia---incluyendo consulta \texttt{LAG} en DuckDB para extracción de características y evaluación de la función logística---registró un tiempo promedio de \SI{0.00042}{\ms} por registro.

Esta latencia equivale a una capacidad superior a 2.38 millones de eventos AVL por segundo en un único núcleo de CPU. Esta eficiencia extrema confirma que el marco EWS propuesto escala para ingerir feeds GTFS-RT simultáneos de múltiples agencias metropolitanas sin requerir clústeres cloud de cómputo pesado, en línea con paradigmas Edge-Cloud de baja latencia descritos por He et al. \cite{26}.

\section{Discusión y Trabajo Futuro}

\subsection{Tiempo de Anticipación Operacional}
La generación de una EWS sólo es valiosa si proporciona tiempo suficiente para que operadores humanos o lazos de control automatizados ejecuten acciones de mitigación. Al predecir el colapso en la parada $n$ usando datos hasta la parada $n-1$, el modelo garantiza un tiempo de anticipación mínimo de 1 parada (típicamente 2 a 4 minutos en entornos urbanos densos). Dado que la latencia algorítmica es prácticamente nula, toda esta ventana temporal queda disponible para estrategias de intervención, como retención dinámica (instruir al bus líder a permanecer más tiempo en plataforma) o skip-stop (instruir al bus rezagado a pasar paradas menores) \cite{10}.

\subsection{Limitaciones en Entornos de Tráfico Mixto}
Aunque la ventana $n-3$ filtra porciones significativas de ruido urbano, la validación se realizó en corredores de tráfico mixto (SF Muni). En estos entornos, factores exógenos fuera de control de la agencia---congestión vehicular, estacionamiento ilegal en paradas y demoras impredecibles de semáforos---introducen una varianza estocástica de base que limita el máximo teórico de recall (actualmente alrededor de $\approx 0.77$).

\subsection{Resiliencia y Limitaciones de API}
El despliegue operativo del pipeline DT propuesto está restringido por límites de tasa del proveedor externo de telemetría. La API abierta 511.org impone una cuota estricta de 60 solicitudes por hora. Dado que una actualización de estado completa requiere dos endpoints distintos (\texttt{VehiclePositions} y \texttt{TripUpdates}), el sistema se limita a 30 ciclos completos por hora, estableciendo un intervalo mínimo teórico de sondeo de 120 segundos. Para amortiguar latencia de red transitoria y evitar throttling, el motor de ingestión opera con ciclo de 130 segundos ($T_{poll}=130$ s).

Este muestreo discreto de baja frecuencia introduce error de cuantización estocástica en el sistema físico de tiempo continuo. Sea $p_{loss}$ la probabilidad de pérdida de paquete o timeout de API durante un ciclo de sondeo. La incertidumbre en el tiempo observado de arribo $T_{obs}$ crece proporcionalmente con actualizaciones consecutivas perdidas. Si se pierden $k$ ciclos consecutivos, el error acotado en drift de retraso $D$ y headway observado $h_{obs}$ puede modelarse como $\epsilon \sim \mathcal{U}(0, k \cdot T_{poll})$. El efecto cascada de esta incertidumbre sobre el índice de agrupamiento ($B_i$) y la desviación de headway ($\Delta H$) representa un riesgo severo de falsas alertas EWS.

Para preservar resiliencia del clasificador, el pipeline implementa un protocolo determinístico de marcado e imputación. Cuando $k=1$, coordenadas espaciales faltantes se aproximan por interpolación lineal basada en geometría GTFS shape y el último vector de velocidad conocido. Los valores correspondientes de $D$ y $h_{obs}$ se recalculan con bandera de incertidumbre. Sin embargo, si $k \ge 2$ ($>260$ s de telemetría faltante), la interpolación se considera físicamente inválida por alta varianza de tráfico urbano. En esos casos, la observación se marca como \textit{stale} y la ventana secuencial $(n-3)$ asociada $X_n$ se excluye explícitamente de la cola de inferencia. Este mecanismo fail-safe garantiza que el modelo reciba sólo datos de alta fidelidad, intercambiando deliberadamente una pérdida temporal de recall por la preservación de precisión de 0.95.

\subsection{Transición hacia Sistemas BRT}
El trabajo futuro se enfocará en calibrar y desplegar el marco EWS propuesto en sistemas Bus Rapid Transit (BRT) con carriles exclusivos segregados físicamente. La literatura reciente demuestra que aplicar estrategias de retención dinámica en corredores BRT segregados reduce significativamente penalidades por agrupamiento \cite{10,29}.

Al migrar este marco algorítmico a un Gemelo Digital BRT, la varianza exógena de tráfico se aísla matemáticamente. Los principales impulsores de $\Delta H$ pasan a ser estrictamente endógenos: tasas de abordaje (tiempos de permanencia) y comportamiento individual de conducción. Hipotetizamos que aplicar este mismo modelo secuencial en topologías de carril exclusivo elevará el F1-Score por encima de 0.90, generando señal suficientemente limpia \cite{30} para activar controles automáticos de retención en lazo cerrado \cite{10} sin validación humana en línea.

\section{Conclusión}
Este artículo estudia la predicción de agrupamiento de buses en entornos de tránsito de tráfico mixto mediante un marco secuencial de Gemelo Digital. Encontramos que una ventana espacio-temporal multi-parada ($n-3$), optimizada con Regresión Logística regularizada L2, filtra eficazmente ruido urbano y logra precisión de transición de 0.95 (IC 95\%: [0.93, 0.97]) y F1-Score corregido de 0.8234 (IC 95\%: [0.80, 0.84]). La evidencia es consistente con la hipótesis de que la tasa de colapso del headway ($\Delta H$) es el precursor dominante singular de degradación sistémica, aunque la varianza estocástica exógena de corredores de tráfico mixto limita la inferencia sobre predicción absoluta de anomalías. Estos hallazgos informan la gestión proactiva del tránsito y el monitoreo Edge-Cloud al proveer señales de alerta temprana computacionalmente eficientes (sub-milisegundo).

\section*{Reproducibilidad}
Para garantizar transparencia metodológica y apoyar investigación futura, los scripts completos de ingestión DuckDB, pipelines de transformación y matrices de configuración de hiperparámetros utilizadas en este estudio están disponibles como código abierto en: \url{https://github.com/gentleman-programming/pp-pipe-dt}.

\begin{thebibliography}{30}

	\bibitem{1} C. F. Daganzo, ``A headway-based approach to eliminate bus bunching: Systematic analysis and comparisons,'' \textit{Transportation Research Part B: Methodological}, vol. 43, no. 10, pp. 913--921, 2009. \url{https://doi.org/10.1016/j.trb.2009.04.002}

	\bibitem{2} J. J. Bartholdi and D. D. Eisenstein, ``A self-coordinating bus route to resist bus bunching,'' \textit{IEEE Transactions on Intelligent Transportation Systems}, vol. 13, no. 2, pp. 489--498, 2012. \url{https://doi.org/10.1109/TITS.2011.2174232}

	\bibitem{3} P. Wang et al., ``Modeling and Prediction of Bus Operation States for Bunching Analysis,'' \textit{IEEE Access}, vol. 8, pp. 10256--10268, 2020. \url{https://doi.org/10.1109/ACCESS.2020.2991708}

	\bibitem{4} H. Jo, ``Navigating the bus bunching problem in urban environments with GTFS-Realtime data,'' \textit{Transportation Research Part C: Emerging Technologies}, vol. 145, 2023.

	\bibitem{5} C. Botín-Sanabria et al., ``Digital-twin–driven urban lifecycle paradigm: evidence from the Singapore–Nanjing eco Hi-Tech Island,'' \textit{Frontiers in Sustainable Cities}, 2026. \url{https://doi.org/10.3389/frsc.2026.1733281}

	\bibitem{6} P. Manandhar et al., ``Leveraging Digital Twin Technology for Sustainable and Efficient Public Transportation,'' \textit{Applied Sciences}, vol. 15, no. 6, 2942, 2025. \url{https://www.mdpi.com/2076-3417/15/6/2942}

	\bibitem{7} L. Hu et al., ``From Raw GPS to GTFS: A Real-World Open Dataset for Bus Travel Time Prediction,'' \textit{Data}, vol. 10, no. 8, 119, 2025. \url{https://doi.org/10.3390/data10080119}

	\bibitem{8} S. V. Kumar and L. Vanajakshi, ``Evaluation of early warning metrics for headway degradation in BRT corridors,'' \textit{Transportation Letters}, vol. 16, no. 8, pp. 889--899, 2024.

	\bibitem{9} J. Zhao, A. Bhowmick, and V. Corcoran, ``Measuring transit performance using AVL and GTFS data,'' \textit{Public Transport}, vol. 12, no. 1, pp. 1--25, 2020.

	\bibitem{10} C. M. Olvera-Toscano et al., ``Holding times to maintain quasi-regular headways and reduce real-time bus bunching,'' \textit{Public Transport}, 2023. \url{https://pmc.ncbi.nlm.nih.gov/articles/PMC10188328/}

	\bibitem{11} A. Aemmer et al., ``Measurement and classification of transit delays using GTFS-RT data,'' \textit{Journal of Public Transportation}, vol. 24, 2022. \url{https://doi.org/10.1007/s12469-022-00291-7}

	\bibitem{12} L. Ge, M. Sarhani, S. Voß, and L. Xie, ``Review of transit data sources: Potentials, challenges and complementarity,'' \textit{Sustainability}, vol. 13, no. 20, 11450, 2021.

	\bibitem{13} A. Graser, ``Analyzing GTFS Realtime data for public transport insights using column-oriented databases,'' \textit{Data Engineering Letters}, vol. 14, pp. 45--56, 2024.

	\bibitem{14} M. Raasveldt and H. Mühleisen, ``DuckDB: an Embeddable Analytical Database,'' in \textit{Proceedings of the 2019 International Conference on Management of Data (SIGMOD)}, 2019, pp. 1981--1984. \url{https://doi.org/10.1145/3299869.3320212}

	\bibitem{15} E. Boudabbous et al., ``Scalable Transit Delay Prediction at City Scale: A Systematic Approach with Multi-Resolution Feature Engineering and Deep Learning,'' \textit{arXiv preprint arXiv:2601.18521}, Jan. 2026. \url{https://arxiv.org/abs/2601.18521}

	\bibitem{16} E. Al-Fares et al., ``Handling class imbalance in stochastic transportation anomaly prediction,'' \textit{Expert Systems with Applications}, vol. 238, 121742, 2024.

	\bibitem{17} H. Chen and Y. Cheng, ``Travel Mode Choice Prediction Using Imbalanced Machine Learning,'' \textit{IEEE Transactions on Intelligent Transportation Systems}, 2023. \url{https://doi.org/10.1109/TITS.2023.3237681}

	\bibitem{18} L. Li, D. Wen, and N. N. Zheng, ``Imbalanced data learning for rare transit anomaly detection,'' \textit{Pattern Recognition}, vol. 122, 2022.

	\bibitem{19} S. Jeong et al., ``BAT-transformer: Prediction of bus arrival time with transformer encoder,'' \textit{IEEE Transactions on Intelligent Transportation Systems}, vol. 24, no. 1, pp. 882--893, 2023.

	\bibitem{20} T. Yin, Y. Li, and J. Zhao, ``Deep learning vs. logistic regression for transit reliability prediction: A comparative study,'' \textit{Transportation Research Part C}, vol. 114, pp. 241--255, 2020.

	\bibitem{21} Y. Yang, J. Cheng, and Y. Liu, ``An overview of solutions to the bus bunching problem in urban bus systems,'' \textit{Frontiers of Engineering Management}, vol. 11, no. 4, p. 661, 2024. \url{https://doi.org/10.1007/s42524-024-0297-1}

	\bibitem{22} Y. Lu et al., ``Digital Twin-driven smart transportation: A review,'' \textit{IEEE Transactions on Intelligent Transportation Systems}, vol. 21, no. 8, pp. 3233--3247, 2020. \url{https://doi.org/10.1109/TITS.2020.2988332}

	\bibitem{23} M. S. Hossain and M. R. Karim, ``Digital Twin framework for proactive urban traffic management: A review,'' \textit{IEEE Access}, vol. 12, pp. 24510--24528, 2024.

	\bibitem{24} A. Kuncheria, ``Simulating the future: The role of Digital Twins in transportation planning,'' \textit{Institute of Transportation Studies Berkeley}, 2025.

	\bibitem{25} M. Al-Azzawi et al., ``Architecting Digital Twins for Intelligent Transportation Systems,'' \textit{arXiv preprint arXiv:2502.17646}, Feb. 2025. \url{https://arxiv.org/abs/2502.17646}

	\bibitem{26} Y. He et al., ``Edge computing architectures for low-latency transit Digital Twins,'' \textit{IEEE Internet of Things Journal}, vol. 11, no. 4, pp. 5812--5825, 2024.

	\bibitem{27} K. Chen, ``Optimizing high-frequency transit operations via Edge-Cloud Digital Twin architectures,'' \textit{Journal of Advanced Transportation}, vol. 2025, 2025.

	\bibitem{28} Bureau of Transportation Statistics, ``Assessing GTFS Accuracy,'' \textit{Mineta Transportation Institute}, 2025. \url{https://rosap.ntl.bts.gov/view/dot/77190}

	\bibitem{29} Y. Zhang, ``Single Agent Robust Deep Reinforcement Learning for Bus Fleet Control,'' \textit{Transportation Safety and Environment}, tdag005, 2026. \url{https://doi.org/10.1093/tse/tdag005}

	\bibitem{30} R. Wang, J. Zhang, and C. Liu, ``Spatio-temporal graph neural networks for short-term bus bunching prediction,'' \textit{Transportation Research Part C: Emerging Technologies}, vol. 158, 104432, 2024.

\end{thebibliography}

\end{document}
